\section[Conclusão]{Conclusão}

Minha experiência como \ac{ages} I foi muito importante no meu crescimento pessoal e profissional. Foi uma oportunidade de vivenciar de perto a dinâmica de trabalho em equipe multidisciplinar, proporcionando uma visão mais realista do que encontrarei no mercado de trabalho. Acompanhei de perto o processo de desenvolvimento de um software a partir de um código existente, o que foi muito enriquecedor.

Fiquei impressionado com o engajamento e compromisso da equipe com o projeto, especialmente em relação à entrega ao \textit{stakeholder}. Essa dedicação nos permitiu estabelecer uma maior coesão entre os membros da equipe. Além disso, tive a oportunidade de conhecer e interagir com colegas de curso que não conhecia antes, criando laços acadêmicos que levo comigo para os próximos semestres.

A experiência de desenvolver um projeto com entregas iterativas e incrementais, com o compromisso de entregar um sistema funcional ao final, foi muito interessante. Embora na Agência não sejamos ``obrigados'' a entregar um sistema completo aos clientes no final do semestre, o fato de construir algo que será utilizado por outras pessoas foi uma motivação poderosa para dar o meu melhor.

Enfrentamos alguns problemas de comunicação dentro da equipe, o que resultou em atrasos na entrega de algumas funcionalidades. O principal obstáculo foi a definição tardia da interface do \textit{backend}. Um aspecto negativo que observei foi a falta de testes de unidade e integração, algo que considero essencial no desenvolvimento de software.

No que se refere à relação entre as disciplinas cursadas e a \ac{ages}, acredito que, por ter iniciado a minha caminhada na Agência mais tarde do que esperava, consegui contribuir de forma mais efetiva para o projeto, mesmo atuando como \ac{ages} I. As disciplinas que mais me ajudaram nesse sentido foram Fundamentos de \ac{ihc} e Gerenciamento de Configuração de Software. A primeira me deu uma base importante para compreender as decisões de interface e usabilidade do aplicativo, auxiliando no desenvolvimento das telas do \textit{frontend}. Já a segunda foi fundamental para que eu lidasse com o controle de versão e o fluxo de trabalho com Git, facilitando a integração do meu código ao repositório e a colaboração com o restante da equipe.

No geral, acredito que este semestre me proporcionou um crescimento pessoal significativo. Melhorei minhas habilidades sociais ao interagir com meus colegas de equipe e participar de reuniões com o \textit{stakeholder}. Além disso, expandi meu conhecimento técnico ao desenvolver habilidades em \textit{backend} e \textit{frontend}, o que me abre mais oportunidades no mercado de trabalho.